% !TeX root = ../动手学强化学习笔记.tex
% 由主文档通过 \input 载入；请编译主文档。
\setcounter{chapter}{0} % 保留原书章号。
\chapter{符号索引与常见概念辨析}
\label{app:rl-notation}

\section{环境、策略与价值}

$\mathcal S$ 为状态空间，$\mathcal A(s)$ 为状态 $s$ 的合法动作集；$S_t,O_t,A_t,R_{t+1}$ 分别为状态、观测、动作和随后奖励，$H_t$ 为可用历史。$p(s',r\mid s,a)$ 是联合转移奖励核，$P(s'\mid s,a)$ 为边缘转移概率，$r(s,a)$ 为期望奖励，$\mu$ 为初始状态分布，$\Omega$ 为观测核。$\pi$ 表示随机策略，$\mu_\theta(s)$ 表示确定性演员；带参数与自变量的记法将后者与初始分布区分。行为策略记为 $b(a\mid s)$，POMDP 的信念记为 $b_t(s)$，高层选项策略记为 $\nu(o\mid s)$。

$G_t$ 是从时刻 $t$ 开始的回报，$\gamma$ 为折扣因子；$v_\pi,q_\pi,A_\pi$ 分别为真实状态价值、动作价值和优势。$V,Q$ 为表格估计，$V_w,Q_w$ 为参数近似，$\theta$ 为策略参数、$w$ 为价值参数，上标减号表示目标参数。$\alpha$ 为学习率，$\lambda$ 在多步与优势估计中表示迹衰减参数，在约束决策章中表示拉格朗日乘子；$\eta$ 专用于熵温度。$\rho_t$ 为动作概率比，$\delta_t$ 为 TD 残差，$e_t$ 为资格迹，$d_{\mu,\pi}$ 为归一化折扣状态访问分布。

\section{容易混淆的概念}

奖励是单步反馈，回报是未来奖励的累计，价值是回报的条件期望，优势是某动作相对于策略平均行为的价值差。状态承载预测未来所需信息，观测是实际获得的信息；使用观测作为状态需要充分性依据。策略规定如何行动，算法规定如何从数据更新策略或其辅助估计，MDP 则规定问题本身。

规划使用模型计算行为后果，学习利用数据修正估计；同策略与离策略描述数据生成策略和目标策略的关系，在线与离线描述能否继续获取交互。自举表示用估计构造目标，不等于模型规划；无模型算法也可以自举。无偏性是估计量的平均性质，一致性是样本量趋于无穷时的极限性质，收敛速度和有限预算表现又是另外的问题。

\section{方法比较的统一维度}

DP 用已知模型计算期望备份；MC 用完整回报采样且不自举；TD 用一步或多步样本加后继估计；Dyna 将真实更新与模型模拟结合。DQN 以网络近似动作价值；策略梯度直接调节分布参数；演员--评论家用学习型评价信号支持策略梯度。PPO 在新采样批次附近限制替代目标的更新收益，TD3 和 SAC 则大量使用历史回放；SAC 还改变为最大熵目标。比较方法时应同时说明数据、目标、表示、成本和假设，而不是只按名称排序。





